\section*{Chapitre \ref{ch:qm:games-auctions-and-information} --- Jeux, enchères et information}

\begin{solution}{exo:qm:games-auctions-and-information:1}
$b(v) = \frac{n-1}nv = 0.9 \times 0.7 = 0.63$~; recette espérée $\frac{n-1}{n+1} = \frac9{11} = 0.818$.
\end{solution}

\begin{solution}{exo:qm:games-auctions-and-information:2}
Chaque joueur joue le mélange $(\frac12, \frac12)$~: tout autre mélange offre à l'adversaire une réponse pure optimale qui gagne en moyenne. La valeur est 0.
\end{solution}

\begin{solution}{exo:qm:games-auctions-and-information:3}
$H = \log_26 = 2.585$ bits. La parité est une fonction du dé, donc $I = H(\text{parité}) - H(\text{parité} \mid \text{dé}) = 1 - 0 = 1$ bit.
\end{solution}

\begin{solution}{exo:qm:games-auctions-and-information:4}
$\int_0^vF^{n-1} = \int_0^vx^{2n-2}dx = v^{2n-1}/(2n-1)$, et en divisant par $F(v)^{n-1} = v^{2n-2}$ on obtient $v/(2n-1)$, donc $b(v) = v\frac{2n-2}{2n-1}$~: avec cinq enchérisseurs $\frac89v$, soit une sous-enchère plus faible que pour des valeurs uniformes
car les valeurs des adversaires sont plus souvent élevées. La formule générale, évaluée numériquement, donne 0.2667, 0.5333 et 0.8000 pour 0.3, 0.6 et 0.9.
\end{solution}

\begin{solution}{exo:qm:games-auctions-and-information:5}
$r - (1 - r) = 0.2$, donc $r^* = 0.6$.
\end{solution}

\begin{solution}{exo:qm:games-auctions-and-information:6}
Avec $k = 1$, l'enchère vaut $E[Y \mid Y < v]$ où $Y$ est la plus haute des $n - 1$ autres valeurs, dont la fonction de répartition est $F^{n-1}$~; $E[Y \mid Y < v] = v - \int_0^vF^{n-1}/F(v)^{n-1}$ par intégration par parties, ce qui donne
l'\omterm{def:qm:games-auctions-and-information:auctions}{enchère au premier prix}.
\end{solution}

\begin{solution}{exo:qm:games-auctions-and-information:7}
\texttt{equivalence()}~: 0.6665 (premier prix), 0.6665 (second prix), 0.6715 (anglaise, pas de cotation de 0.01), 0.6614 (hollandaise, pas de cotation de 0.01), contre $\frac23$~; avec le \omterm{def:qm:games-auctions-and-information:open}{prix de réserve} 0.5, l'\omterm{def:qm:games-auctions-and-information:auctions}{enchère au second prix} donne 0.6720 contre la
forme fermée 0.6719, avec une vente dans 96.9~\% des cas.
\end{solution}

\begin{solution}{exo:qm:games-auctions-and-information:8}
Gagner chaque enchère à la fair value signifie que les valorisations de la société étaient chaque fois les plus élevées, justement le moment où une valorisation bruitée surestime la vérité~: la malédiction du vainqueur. Un taux de victoire proche de 100~\%
face à quatre concurrents signale que les enchères sont trop élevées, et non que le modèle est bon~; le bon contrôle est le profit réalisé sur les transactions remportées.
\end{solution}

\begin{solution}{pb:qm:games-auctions-and-information:1}
\textbf{1.} $\frac45$ de sa valeur~: l'équilibre arbitre la probabilité de gagner contre la marge, et récupère l'écart espéré avec la valeur suivante.
\textbf{2.} L'enchère détermine seulement si l'on gagne, pas ce que l'on paie~; s'en écarter ne permet de gagner qu'à perte, ou de ne perdre qu'un profit.
\textbf{3.} Le profit espéré culmine pour une enchère de 0.64, l'enchère d'équilibre~: la stratégie est une meilleure réponse.
\textbf{4.} $\frac23$ dans chaque format~; un million d'enchères donnent 0.6665 (premier prix) et 0.6665 (second prix).
\textbf{5.} L'horloge anglaise s'arrête au premier pas de cotation au-dessus de la seconde valeur (un demi-pas de cotation trop haut en moyenne), l'horloge hollandaise au premier pas de cotation sous l'enchère (un demi-pas de cotation trop bas)~: 0.6715 et 0.6614.
\textbf{6.} $r^* = 0.5$~: de $\frac13$ à $\frac5{12}$ (25~\%) avec deux enchérisseurs, l'objet reste invendu un quart du temps~; de 0.6667 à 0.6719 (0.8~\%) avec cinq, il reste invendu 3.1~\% du temps.
\textbf{7.} $b(v) = v(3 - 2.4v)/(4 - 3v)$~: 0.36 pour 0.5, 0.54 pour 0.8, 0.60 pour 1.
\textbf{8.} 1 en espérance dans les deux cas~; la simulation donne 0.9994 (à prix uniforme) et 0.9999 (au prix offert).
\textbf{9.} Le format à prix uniforme disperse les enchères environ deux fois plus (écart entre enchères gagnantes 0.167 contre 0.087) et sa recette est plus de deux fois plus volatile (0.378 contre 0.160)~: enchères plus agressives et plus dispersées
ainsi qu'une volatilité plus forte d'une enchère à l'autre, comme le rapportait l'étude.
\textbf{10.} Avec une demande unitaire, des enchérisseurs symétriques neutres au risque et des valeurs indépendantes, l'équivalence des recettes rend les formats égaux en moyenne~; la réponse dépend de la demande multi-unités, de l'aversion au risque, de l'asymétrie et
de l'information, ce qui explique que le Trésor américain ait mené une expérience.
\textbf{11.} Enchérir son propre signal perd 0.067 par victoire (la surestimation espérée du plus haut des cinq signaux), dans 97~\% des victoires~; une sous-enchère de $0.1 \times \frac46$ ramène le profit à zéro.
\textbf{12.} 1.846 bits~; $-0.165$ bits par course (la commission dépasse le savoir du parieur).
\textbf{13.} 0.363 bits, ce qui porte le taux de doublement à $+0.198$~; la simulation donne un gain de $0.363 \pm 0.002$.
\textbf{14.} Ils multiplient les deux taux de doublement par un même facteur dans le logarithme, qui s'annule dans la différence.
\textbf{15.} Rien~: avec quatre chevaux, un tuyau juste avec probabilité $\frac14$, et uniformément faux sinon, est indépendant du gagnant.
\textbf{16.} L'aversion au risque des enchérisseurs, ou la réduction de la demande multi-unités sous l'\omterm{def:qm:games-auctions-and-information:pab}{enchère au prix offert}, une participation plus large qu'attire le prix uniforme, et une malédiction du vainqueur plus faible lorsque chaque gagnant paie le même prix.
\textbf{17.} Dès que les autres enchérisseurs s'appuient sur la même information~: il faut réduire son enchère de la surestimation espérée de l'estimation gagnante, d'autant plus que les concurrents sont nombreux et que les valorisations sont bruitées.
\textbf{18.} Moins que la croissance qu'il apporte~: au plus son \omterm{def:qm:games-auctions-and-information:entropy}{information mutuelle} avec le résultat, en taux de doublement, net de ce que le tuyau coûte par course.
\textbf{19.} Résultat nommé~: \emph{le basculement du Trésor}~: dans le modèle symétrique à valeurs privées, les deux formats rapportent la même recette espérée (1 pour deux unités et cinq dealers~; $\frac23$ dans le cas à une seule unité), chaque dealer
réduit son \omterm{def:qm:games-auctions-and-information:auctions}{enchère au premier prix} d'un cinquième, et le \omterm{def:qm:games-auctions-and-information:open}{prix de réserve} optimal de 0.5 ajoute 0.8~\% avec cinq dealers (25~\% avec deux).
\textbf{20.} Dans ses conditions, le format n'a pas d'importance pour la recette espérée~; ce qui compte, c'est qui participe, ce qu'ils savent et le \omterm{def:qm:games-auctions-and-information:open}{prix de réserve}.
\end{solution}

\begin{solution}{iq:qm:games-auctions-and-information:1}
Enchérir $v/2$~: face à un adversaire qui enchérit $w/2$, enchérir $b$ gagne avec probabilité $2b$ et rapporte $(v - b)2b$, maximal pour $b = v/2$.
\iqlookfor{Le compromis et la réponse $v/2$.}
\end{solution}

\begin{solution}{iq:qm:games-auctions-and-information:2}
Dans une \omterm{def:qm:games-auctions-and-information:auctions}{enchère à valeur commune}, le gagnant est celui qui a l'estimation la plus optimiste, si bien que, conditionnellement à la victoire, l'estimation est biaisée à la hausse. Il faut enchérir comme si son estimation était la plus haute~: retrancher la
surestimation espérée du maximum des estimations des concurrents, qui croît avec leur nombre et leur bruit.
\iqlookfor{\omterm{def:qm:floating-point-and-numerical-linear-algebra:cond}{Conditionnement} sur la victoire.}
\end{solution}

\begin{solution}{iq:qm:games-auctions-and-information:3}
Toute enchère qui attribue le bien à la plus haute valeur et laisse un surplus nul au type le plus bas rapporte la même recette espérée, avec des valeurs privées indépendantes et des enchérisseurs symétriques neutres au risque. Elle ne tient plus en cas de
l'aversion au risque (le premier prix rapporte davantage), des valeurs affiliées ou communes (l'anglaise rapporte davantage), des enchérisseurs asymétriques, des contraintes budgétaires et une demande multi-unités.
\iqlookfor{Hypothèses et deux cas d'échec.}
\end{solution}

\begin{solution}{iq:qm:games-auctions-and-information:4}
Son \omterm{def:qm:games-auctions-and-information:entropy}{information mutuelle} avec le résultat, en taux de doublement~: la croissance logarithmique optimale augmente exactement de $I(X; Y)$ par pari, indépendamment de la cote.
\iqlookfor{\omterm{def:qm:games-auctions-and-information:entropy}{Information mutuelle}.}
\end{solution}

\begin{solution}{iq:qm:games-auctions-and-information:5}
Pas de point-selle (maximin $-1$, minimax 1). Le joueur en ligne joue la première ligne avec probabilité $x = \frac37$, le joueur en colonne joue la première colonne avec probabilité $y = \frac27$~; la valeur est $\frac17$.
\iqlookfor{Conditions d'indifférence.}
\end{solution}

\begin{solution}{iq:qm:games-auctions-and-information:6}
Dans une \omterm{def:qm:games-auctions-and-information:auctions}{enchère à valeurs privées} au premier prix, davantage de concurrents signifie moins de sous-enchère (enchère plus proche de la valeur)~; avec des valeurs communes (l'essentiel de la valeur du portefeuille est le même pour tous les dealers), davantage de concurrents signifie une malédiction du
vainqueur plus forte et davantage de sous-enchère. Une demande de cotation (RFQ) de portefeuille comporte une large composante commune, si bien que le second effet domine en général.
\iqlookfor{Les deux effets opposés.}
\end{solution}
